% Noether Paper 11 controlled-generic zh-Hant producer translation.
% Current-German interval SHA-256:
% DE540AFFE512CF39871019A946F7CEBA9FFF04ADFB54FA1D531E5A40ADA5CBA1.
% Inherited Chinese content witness SHA-256:
% D0737F8288235A705A1E2371BCD9F535CB5375212BBEC983324B7077014A23DD.
% Controlled-generic script transport only; every independent check pending.
% Controlled generic Traditional script only; not zh-Hant-TW/HK/MO prose.
% Hans wording remains the lexical base; independent Hant checking is pending.
% Compilation is mechanical production only; no source, semantic, formula-content,
% terminology, visual, regional, publication, archive, or certification check is claimed.
\documentclass[11pt]{article}
\usepackage[a4paper,margin=2.4cm]{geometry}
\usepackage{fontspec}
\usepackage{xeCJK}
\usepackage{amsmath,amssymb,mathtools,array}
\setmainfont{Noto Serif}
\setCJKmainfont{Microsoft JhengHei}
\setCJKsansfont{Microsoft JhengHei}
\XeTeXlinebreaklocale "zh"
\XeTeXlinebreakskip=0pt plus 1pt
\setlength{\parindent}{2em}
\setlength{\parskip}{0.35em}
\setlength{\emergencystretch}{8em}
\allowdisplaybreaks
\raggedbottom
\providecommand{\glqq}{“}
\providecommand{\grqq}{”}
\providecommand{\frK}{\mathfrak K}
\providecommand{\srcfn}[2]{%
  \begingroup
  \renewcommand{\thefootnote}{#1}%
  \footnote{#2}%
  \addtocounter{footnote}{-1}%
  \endgroup}
\providecommand{\srcfnmark}[1]{%
  \begingroup
  \renewcommand{\thefootnote}{#1}%
  \footnotemark%
  \addtocounter{footnote}{-1}%
  \endgroup}
\providecommand{\srcfntext}[2]{%
  \begingroup
  \renewcommand{\thefootnote}{#1}%
  \footnotetext{#2}%
  \endgroup}
\providecommand{\srcnumdisplay}[3][3.4em]{%
  \par\smallskip\noindent
  \begin{minipage}[t]{#1}\vspace*{0pt}#2\end{minipage}%
  \begin{minipage}[t]{\dimexpr\linewidth-#1\relax}\vspace*{0pt}\centering
  \(\displaystyle #3\)
  \end{minipage}\par\smallskip}
\begin{document}
\setcounter{footnote}{0}
\section*{11. 具有給定群的方程。}
\begin{center}
{\Large\bfseries 11. 具有給定群的方程。}\par
\vspace{0.75em}
作者\par
\vspace{0.75em}
哥廷根的 Emmy Noether，\par
\vspace{1.5em}
Math. Ann. 78 (1918)，第 221--229 頁
\end{center}

\vspace{4em}

構造具有給定群的方程這一問題，可以從兩個方向入手；可簡要地將其稱為刻畫根的\glqq{}無理\grqq{}方向和刻畫係數的\glqq{}有理\grqq{}方向。

在以函數論和算術方法研究的\glqq{}無理\grqq{}方向中，有 Kronecker 定理，即有理數域上的所有阿貝爾域都是分圓域；此外還有關於二次數域的相對阿貝爾域的相應定理。這些定理一方面給出了所有阿貝爾群相對於這些特殊有理性域的可實現性；另一方面又給出了方程的全體，同時使人對由此定義的數域結構獲得更深入的認識。然而，每一個這樣的定理都侷限於一個特定的有理性域，而每一個新的有理性域都需要一種全新的處理方法。

以代數方法研究的\glqq{}有理\grqq{}方向以 Hilbert 不可約性定理為基礎；根據該定理，在係數表示中可以只考慮參數表示。例如，相對於任意有理性域，存在任意多個具有交錯群的方程，就屬於這一方向。在這裡，自然缺少上文所述的對方程所定義之域的結構認識；不過，這個\glqq{}有理\grqq{}方向的優點在於，它同時對\emph{所有有理性域}證明了群的可實現性；但迄今已知的參數表示一般並不給出方程的全體。\srcfn{*)}{對於可解方程，可以用根的參數表示代替係數的參數表示。最近，F. Mertens 對某些八次群給出了最一般的根表示：\glqq{}具有四元數群的八次方程\grqq{}，Sitzb. d. Ak. d. Wiss., Wien, Abt. IIa, Bd. 125, S. 735。據我從這篇札記得知，在此以前 G. Bucht 已經對三次、四次方程以及具有上述群的八次方程中總能構造的標準形，給出了最一般的根表達式：\glqq{}關於若干八次代數域\grqq{}，Arkiv for Math., Astron. och Physik, Bd. 6, Nr. 30。[校樣時補記。]} 

下文是對\glqq{}有理\grqq{}方向的一項貢獻；這裡將給出\emph{充分條件}，以保證\emph{存在這樣的參數表示：對於任意有理性域，它們都能給出具有給定群的方程全體}。這些條件是：屬於該群的不變量域——Lagrange 屬域——能夠由該域中的\emph{獨立}函數作有理表示，即具有一個\glqq{}最小基\grqq{}；換言之，它同構於 $n$ 個不定元的所有有理函數所成的域。如 § 2 中還將說明的，例如三次和四次方程中出現的所有群都滿足這一條件。關於這些三次和四次方程的實際構造和更詳細討論，可見 F. Seidelmann 的學位論文，\srcfn{*)}{任意有理性域上帶約束的三次方程與雙二次方程全體。埃爾朗根，1916 年。}該論文的摘錄緊接在本通信之後。

正如論文\glqq{}有理函數的域與系統\grqq{}（Math. Ann. 76）中已經提到的，Lagrange 屬域的最小基問題，是 E. Fischer 在與我的交流中獨立於上述聯繫而提出的，併成為本研究的出發點。\srcfn{**)}{參見關於\glqq{}有理函數域\grqq{}的預告通信第 2 部分，Jahresb. d. d. Math. Vereinig., Bd. 22, 1913。}

\begin{center}
§ 1.\\[0.75em]
\textbf{参数表示的充分条件。}
\end{center}

1. 将具有给定群的方程之构造化归为参数表示的可能性，来自 Hilbert 不可约性定理；该定理陈述如下：设方程
\[
\begin{array}{@{}l@{\quad}l@{}}
(1)& \displaystyle f(x)=x^n+F_1(\lambda_1\cdots\lambda_r)x^{n-1}+\cdots+F_n(\lambda_1\cdots\lambda_r)=0
\end{array}
\]
——其中 $F_i(\lambda_1\cdots\lambda_r)$ 表示獨立參數 $\lambda_1\cdots\lambda_r$ 的有理函數，其係數來自一個數域\srcfn{***)}{也可以用一般的有理性域代替數域。}$\Omega$——相對於由 $\lambda_1\cdots\lambda_r$ 的所有有理函數組成且係數來自 $\Omega$ 的域 $\Omega(\lambda_1\cdots\lambda_r)$，具有群 $\Gamma$；那麼，可以用有理數以任意多種方式替換參數 $\lambda$，使所得數值方程
\[
\begin{array}{@{}l@{\quad}l@{}}
(2)& \displaystyle g(x)=x^n+a_1x^{n-1}+\cdots+a_n=0
\end{array}
\]
相對於 $\Omega$ 恰好具有群 $\Gamma$。這樣的 $g(x)$ 更確切地記作 $g_\Gamma$。

現在要問，能否給出\emph{這樣一種參數表示} (1)，使得\emph{當 $\lambda_1\cdots\lambda_r$ 遍歷 $\Omega$ 中的所有量時，所得數值方程} (2) \emph{遍歷所有 $g_\Gamma$}——但排除 $\Omega$ 中那些導致群縮小的量。\srcfn{*)}{這裡還應包括 $g(x)=0$ 出現重根的情形。}換言之，非線性方程組
\[
\begin{array}{@{}l@{\quad}l@{}}
(3)& \displaystyle F_1(\lambda_1\cdots\lambda_r)=a_1;\quad \cdots;\quad F_n(\lambda_1\cdots\lambda_r)=a_n
\end{array}
\]
對於每一個對應於某個 $g_\Gamma$ 的值組 $a_1\cdots a_n$，都應當在 $\Omega$ 中的量 $\lambda_1\cdots\lambda_r$ 中有解。由於方程組 (3) 是非線性的，若要求通過一個參數表示得到所有 $g_\Gamma$，就必須從一開始作出限制。事實上，在這樣的非線性方程組中，總可能出現滿足某個代數關係 $H(a_1\cdots a_n)=0$ 的奇異值組 $a_1\cdots a_n$，而對這些值組卻\emph{不存在}解；\srcfn{**)}{這一點將於近期在一篇關於\glqq{}非線性方程組中的擇一定理\grqq{}的論文中作一般性闡述。}於是參數表示 (1) 失效，而必須補充另一種表示。不過在當前情形下，如下文將說明的，這些 $a_1\cdots a_n$ 的奇異值可以得到簡單的刻畫。

2. 為了獲得這樣的參數表示，我們進一步要求，$\lambda_i$ \emph{恆等地是關於那些被視為不定元的根的、屬於該群的自然無理量}。其含義如下：若在 (1) 中，用適當選取的、不定元 $x_1\cdots x_n$ 的\emph{有理}函數 $\varphi_1(x)\cdots\varphi_r(x)$——其係數來自 $\Omega$——替換參數 $\lambda_1\cdots\lambda_r$，則 (1) 應由此變為
\[
\begin{array}{@{}l@{\quad}l@{}}
(4)& \displaystyle h(x)=x^n-\sigma_1(x)x^{n-1}+\sigma_2(x)x^{n-2}\cdots\pm\sigma_n(x)
\end{array}
\]
——其中 $\sigma_1(x)\cdots\sigma_n(x)$ 表示 $x_1\cdots x_n$ 的初等對稱函數——並且 $h(x)$ 相對於由 $\Omega(\lambda_1\cdots\lambda_r)$ 以這種方式產生的有理性域 $\Omega(\varphi_1(x)\cdots\varphi_r(x))$ 應具有群 $\Gamma$；由於參數相互獨立，$\varphi_1(x)\cdots\varphi_r(x)$ 也應當是代數無關函數。

為了使這一要求精確化，必須闡明 $\Omega(\varphi_1(x)\cdots\varphi_r(x))$ 與群 $\Gamma$ 的關係。為此，我們首先確定最一般的有理性域 $K$，使 $h(x)$ 相對於它成為 $h_\Gamma$。令 $\Omega_\Gamma$ 表示 $x_1\cdots x_n$ 的所有有理數係數有理函數中在 $\Gamma$ 作用下不變者所成的域——它也稱為屬於 $\Gamma$ 的\glqq{}Lagrange 屬域\grqq{}或 $\Gamma$ 的不變量域——並令 $\Omega(\Omega_\Gamma)$ 表示把 $\Omega_\Gamma$ 添入 $\Omega$ 所得到的相應域。因此，$\Omega_\Gamma$ 是由 $x_1\cdots x_n$ 的所有有理數係數有理函數組成的域 $\mathfrak R(x_1\cdots x_n)$ 的一個子域。根據群的定義，最一般的域 $K$ 可由\emph{任何滿足以下條件的域}給出：它包含 $\Omega_\Gamma$ 作為子域，並且它同 $\mathfrak R(x_1\cdots x_n)$ 的交恰好是 $\Omega_\Gamma$。相應地，對於每一個包含 $\Omega$ 的域 $K$，它同 $\Omega(x_1\cdots x_n)$ 的交由 $\Omega(\Omega_\Gamma)$ 給出。特別地，由於按我們的要求 $\Omega(\varphi_1(x)\cdots\varphi_r(x))$ 是這樣一個域 $K$，它同 $\Omega(x_1\cdots x_n)$ 的交便由 $\Omega(\Omega_\Gamma)$ 給出；另一方面，由於 $\Omega(\varphi_1(x)\cdots\varphi_r(x))$ 是 $\Omega(x_1\cdots x_n)$ 的一個子域，所以 $\Omega(\varphi_1(x)\cdots\varphi_r(x))$ 恰好等於 $\Omega(\Omega_\Gamma)$。因此，這一要求的含義是：\emph{Lagrange 屬域 $\Omega(\Omega_\Gamma)$ 應當由 $\Omega$ 添入代數無關函數 $\varphi_1(x)\cdots\varphi_r(x)$ 而得到}；這裡必須有 $r=n$，因為 $\Omega_\Gamma$ 包含 $n$ 個代數無關的初等函數，而另一方面，不可能有超過 $n$ 個函數代數無關。我們把 $\varphi_1(x)\cdots\varphi_n(x)$ 稱為 $\Omega(\Omega_\Gamma)$ 的一個\emph{最小基}；也可以說，$\Omega_\Gamma$ 具有一個係數來自 $\Omega$ 的最小基。

3. 現在的關鍵是，以上全部論證可以反過來，並由此確實導出所求參數表示的充分條件。\srcfn{*)}{第 2 節中的要求比第 1 節有所加強，因此這些條件不必是必要條件。}於是，我們假定 $\varphi_1(x)\cdots\varphi_n(x)$ 是 $\Omega(\Omega_\Gamma)$ 的一個最小基，並且 $\Omega$ 是包含 $\varphi_1(x)\cdots\varphi_n(x)$ 的係數的最小域；此時 $\Omega$ 必然是一個代數數域，尤其可以是所有有理數所成的域 $\mathfrak R$。由於 $\sigma_1(x),\ldots,\sigma_n(x)$ 屬於 $\Omega(\Omega_\Gamma)$，便存在關於 $x_1\cdots x_n$ 的恆等式——其中 $F_i$ 表示其自變量的有理函數：
\[
\begin{array}{@{}l@{\quad}l@{}}
(5)& \displaystyle \sigma_1(x)=F_1(\varphi_1(x)\cdots\varphi_n(x));\quad \cdots;\quad \sigma_n(x)=F_n(\varphi_1(x)\cdots\varphi_n(x));
\end{array}
\]
反過來，$\varphi_1(x)\cdots\varphi_n(x)$，因而還有 $\varphi(x)=\mu_1\varphi_1(x)+\cdots+\mu_n\varphi_n(x)$，都通過下述相對於 $\Omega(\sigma_1(x)\cdots\sigma_n(x))$ 不可約的方程而代數地依賴於 $\sigma_1(x)\cdots\sigma_n(x)$：
\[
\begin{array}{@{}l@{\quad}l@{}}
(6)& \displaystyle G(\varphi(x);\sigma_1(x)\cdots\sigma_n(x))=0,
\end{array}
\]
該方程仍然是關於 $x$ 的恆等式。

由 (5)，恆等式 (6) 化為
\[
\begin{array}{@{}l@{\quad}l@{}}
(7)& \displaystyle G(\varphi(x);F_1(\varphi_1\cdots\varphi_n);\cdots;F_n(\varphi_1\cdots\varphi_n)) =H(\varphi_1(x)\cdots\varphi_n(x))=0;
\end{array}
\]
由於 $\varphi_1(x)\cdots\varphi_n(x)$ 在代數上無關，(7) 不僅關於
$x$ 恆等成立，而且關於各 $\varphi_i$ 也恆等成立；也就是說，關於
獨立參數 $\lambda$ 有恆等式：
\[
\begin{array}{@{}l@{\quad}l@{}}
(8)& \displaystyle H(\lambda_1\cdots\lambda_n)=G(\mu_1\lambda_1+\cdots+\mu_n\lambda_n;F_1(\lambda_1\cdots\lambda_n);\cdots;F_n(\lambda_1\cdots\lambda_n))=0.
\end{array}
\]
由恆等式 (8) 可以推出，方程
\[
\begin{array}{@{}l@{\quad}l@{}}
(9)& \displaystyle f(x)=x^n-F_1(\lambda_1\cdots\lambda_n)x^{n-1}+F_2(\lambda_1\cdots\lambda_n)x^{n-2}\cdots\pm F_n(\lambda_1\cdots\lambda_n)=0
\end{array}
\]
\emph{相對於} $\Omega(\lambda_1\cdots\lambda_n)$ \emph{具有群}
$\Gamma$。事實上，這個恆等式表明，可以有理地給出的
$\lambda=\mu_1\lambda_1+\cdots+\mu_n\lambda_n$ 是一個屬於
$\Gamma$ 的根函數；而由代換 $\lambda_i=\varphi_i(x)$ 可見，不可能
再有一個屬於 $\Gamma$ 的某個子群的函數也能有理地給出；同時，在同一
代換下，$\lambda$ 還與它的所有共軛量互不相同。此外，若 $\Omega^*$
是任意一個包含 $\Omega$ 的數域，則 $f(x)$ \emph{相對於}
$\Omega^*(\lambda_1\cdots\lambda_n)$ \emph{具有群} $\Gamma$；
因為由第 2 節可知，在作代換 $\lambda_i=\varphi_i(x)$ 時，即使添入
$\Omega^*$，群也不會約化。

現在還須證明，在第 1 節所修正的意義下，$f(x)$ 確實遍歷所有
$g_\Gamma$。為此要注意，由 (5) 可知，$F_i(\lambda)$ 是關於其自變量
$\lambda$ 的代數無關函數。因此，方程組
\[
\begin{array}{@{}l@{\quad}l@{}}
(10)& \displaystyle F_1(\lambda_1\cdots\lambda_n)=-a_1;\quad F_2(\lambda_1\cdots\lambda_n)=a_2;\quad \cdots;\quad F_n(\lambda_1\cdots\lambda_n)=\pm a_n
\end{array}
\]
對於任意值組 $a_1\cdots a_n$ 都有解，只有下文尚待指出的奇異值組
例外。對於相對於 $\Omega^*$ 的某個 $g_\Gamma$ 所對應的每個值組
$a_1\cdots a_n$，相應的 $\lambda_i$ 作為屬於該群的自然無理量，
\srcfn{*)}{另見下文！} 都是 $\Omega^*$ 中的量。因此，當
$\lambda_1\cdots\lambda_n$ 遍歷所有值組時，$g(x)$ 就遍歷所有方程
（在修正後的意義下）；其中包含對應於 $\Omega^*$ 中各值的全部
$g_\Gamma$。反過來，根據 Hilbert 不可約性定理，$\Omega^*$ 中的
$\lambda$ 值\glqq{}一般地\grqq{}也對應於方程 $g_\Gamma$；所以，
在第 1 節所修正的意義下，參數表示 (9) 確實給出了相對於 $\Omega^*$
的全部 $g_\Gamma$。

現在只剩下確定使參數表示失效的奇異值組。將最小基中的函數拆成分子
和分母，寫為
\[
\varphi_i(x)=\frac{\psi_i(x)}{\chi_i(x)}.
\]
把它們代入恆等式 (5)，設這些恆等式（不作約分）化為：
\[
\begin{array}{@{}l@{\quad}l@{}}
(11)& \displaystyle \sigma_k(x)=\frac{G_k(x)}{H_k(x)}\quad\text{或}\quad H_k(x)\cdot\sigma_k(x)=G_k(x),
\end{array}
\]
於是所有 $H_k(x)$ 的乘積可以被所有分母 $\chi_i(x)$ 的乘積整除。
對於所有滿足
\[
H_1(\alpha)\cdot H_2(\alpha)\cdots H_n(\alpha)\neq0
\]
的根組 $\alpha_1\cdots\alpha_n$，所有分母 $\chi_i(\alpha)$ 也都
不消失，因而可以在 (11) 中除以 $H_k(\alpha)$，得到
\[
\begin{array}{@{}l@{\quad}l@{}}
(12)& \displaystyle \sigma_k(\alpha)=\frac{G_k(\alpha)}{H_k(\alpha)}=F_k(\varphi_1(\alpha);\cdots;\varphi_n(\alpha)),
\end{array}
\]
其中，由於分母不消失，各 $\varphi_i(\alpha)$ 都取有限值；對於方程
$g_\Gamma$，這些值還是 $\Omega^*$ 中的量。只要確定
$\alpha_1\cdots\alpha_n$，使 $a_k=(-1)^k\sigma_k(\alpha)$，(12)
就給出 (10) 的解系。\srcfn{*)}{由此，方程的求解便歸結為求解由獨立
函數組成的方程組：
\[
\varphi_1(\alpha_1\cdots\alpha_n)=\lambda_1;\quad \cdots;\quad \varphi_n(\alpha_1\cdots\alpha_n)=\lambda_n.
\]}

因此，只有滿足
\[
H_1(\alpha)\cdot H_2(\alpha)\cdots H_n(\alpha)=0;
\]
的係數組 $a_k=(-1)^k\sigma_k(\alpha)$ 才可能成為\emph{奇異的}；
對這些係數組，(11) 不再給出上述表示，而變成 $0=0$。這時必須另作
研究，才能判定這樣定義的奇異 $a_1\cdots a_n$ 中是否包含對應於方程
$g_\Gamma$ 的值組，以及它們對應於哪些數域。\srcfn{**)}{例如，正如
F. Seidelmann 在前引處所證明的，這種對應於奇異值組的方程只在三次
和四次方程的交錯群情形下出現，而且只對那些包含三次單位根域的數域
$\Omega^*$ 出現。每種情形都可以給出一個簡單的補充表示。} 綜上可得

\emph{定理：屬於群 $\Gamma$ 的 Lagrange 型域若有一個最小基}
\srcfn{***)}{從一個最小基的存在立刻可以推出任意多個最小基的存在；
它們都能由原來的最小基通過可逆有理變換導出。}
\emph{，其係數屬於 $\Omega$，那麼對於每個包含 $\Omega$ 的數域
$\Omega^*$，相對於 $\Omega^*$、具有預先指定的群 $\Gamma$ 的全部
方程，都可以通過參數表示 (2) 作有理構造——可能例外的是相對於該參數
表示為奇異的方程，但這些方程可以另行給出。參數必須遍歷 $\Omega^*$
中所有不導致群約化的值。若最小基的係數域 $\Omega$ 等於所有有理數
所成的域 $\mathfrak R$，則對於任意數域都可以作這種構造。}

根據 Hilbert 不可約性定理，從參數流形中排除 $\Omega^*$ 內那些導致
群約化的值，並不會降低該流形的維數。因此還有：
\emph{最小基的存在推出，相對於 $\Omega^*$、具有群 $\Gamma$ 的方程
所成流形的維數，等於無特約方程所成流形的維數；特約並不降低這一
維數。}

\begin{center}
§ 2.\\[0.75em]
\textbf{在特殊方程上的应用。}
\end{center}

现在我们证明，在三次和四次方程中出现的所有群都存在最小基；为此，我们首先完全一般地说明，对于 $n^{\text{ten}}$ 次方程，问题可以约化为含有 $n-2$ 个不定元的问题。第一次约化对应于用线性 Tschirnhaus 变换消去次高次项。我们取第一基函数 $\varphi_1(x)=\sigma_1(x)$；并进一步注意到，$\Omega_\Gamma$ 可以从 $x_1\cdots x_n$ 的所有\emph{齐次且容许 $\Gamma$ 的}形式有理地导出。设 $p(x_1\cdots x_n)$ 是这样的一个形式；那么
\[
q(x)=p\left(x_1-\frac{\sigma_1(x)}{n};\cdots;x_n-\frac{\sigma_1(x)}{n}\right)=p\left(x_i-\frac{\sigma_1(x)}{n}\right)=p(y_i)
\]
由於它容許 $\Gamma$，也屬於 $\Omega_\Gamma$，而且有：
\[
p(x)\equiv p\left(x_i-\frac{\sigma_1(x)}{n}\right)\operatorname{mod}\sigma_1(x)
\]
或者：
\[
p(x)=q(x)+\sigma_1(x)\cdot p_1(x)
\]
其中 $p_1(x)$ 屬於 $\Omega_\Gamma$；而且它的次數低於 $p(x)$。繼續這一過程便可看出，$\Omega_\Gamma$ 中的所有齊次形式 -- 從而 $\Omega_\Gamma$ 中的所有有理函數 -- 都可以由 $\varphi_1(x)=\sigma_1(x)$ 以及齊次形式
\[
q(x)=p\left(x_i-\frac{\sigma_1(x)}{n}\right)=p(y_i).
\]
有理地表示。然而，這些組合 $y_i$ 之間存在一個線性關係：$\sum y_i=0$；因此，$q(x)$ \emph{只依賴於 $(n-1)$ 個自變量。}

第二次約化以 $q(x)$ 是其自變量的齊次形式這一事實為基礎。我們取第二基函數
\[
\varphi_2(x)=\frac{\tau_3(x)}{\tau_2(x)},\quad \text{其中}\quad
\tau_3(x)=\sigma_3\left(x_i-\frac{\sigma_1(x)}{n}\right)=\sigma_3(y_i),\quad
\tau_2(x)=\sigma_2\left(x_i-\frac{\sigma_1(x)}{n}\right)=\sigma_2(y_i);
\]
因此，$\varphi_2(x)$ 是關於 $(n-1)$ 個自變量的一次齊次分式函數，這些自變量就是滿足 $\sum y_i=0$ 的 $y_i$。藉助 $\varphi_2(x)$，所有 $q(x)$ 都可以由同樣屬於 $\Omega_\Gamma$ 的函數
\[
r(x)=\frac{q(x)}{\varphi_2(x)^\lambda}=\frac{q(x)\cdot\tau_2(x)^\lambda}{\tau_3(x)^\lambda}=\frac{p(y_i)\cdot\sigma_2(y_i)^\lambda}{\sigma_3(y_i)^\lambda},
\]
有理地表示，其中 $\lambda$ 表示 $q(x)$ 的次數。於是 $r(x)$ 成為零次齊次函數，因此只依賴於 $(n-2)$ 個自變量，即在 $\sum y_i=0$ 條件下的比值 $y_1:y_2:\cdots:y_n$。所以，$\Omega_\Gamma$ 可以由
\[
\varphi_1(x)=\sigma_1(x);\quad \varphi_2(x)=\frac{\tau_3(x)}{\tau_2(x)}
\]
以及\emph{$\Omega_\Gamma$ 中所有依賴於這 $(n-2)$ 個自變量的函數 $r(x)$ 所組成的域}有理地表示，\srcfn{*)}{對於不含次高次項的 $n^{\text{ten}}$ 次方程：
\[
f(x)=x^n+a_2x^{n-2}+a_3x^{n-3}+\cdots+a_n=0
\]
這第二次約化對應於代換 $y=\frac{x\cdot a_2}{a_3}$；當 $a_2\neq0$、$a_3\neq0$ 時，這一代換總是可行的。由此，$f(x)$ 在相差一個因子的意義下變為
\[
g(y)=y^n+\frac{a_2^3}{a_3^2}y^{n-2}+\frac{a_2^3}{a_3^2}y^{n-3}+\frac{a_4\cdot a_2^4}{a_3^4}y^{n-4}+\cdots+a_n\frac{a_2^n}{a_3^n}=0;
\]
也就是說，變為一個只含有 $n-2$ 個參數的方程：
\[
g(y)=y^n+c\cdot y^{n-2}+c\cdot y^{n-3}+c_4y^{n-4}+\cdots+c_n=0.
\]} 這樣便完成了所要求的約化。

對於三次和四次方程，由此特別得到向一個或兩個自變量的函數域的約化。根據 Lüroth 和 Castelnuovo 的定理，這些函數域總是存在最小基。\srcfn{**)}{J. Lüroth: Beweis eines Satzes über rationale Kurven, Math. Ann. 9 (1875); G. Castelnuovo: Sulla razionalità delle involuzioni piane, Math. Ann. 44 (1893). F. Enriques 用一個反例表明，含有兩個以上不定元的域不一定存在最小基，Rend. Acc. Linc., Vol. 21, 21. Jan. 1912.} 具體地說，在一個不定元的情形下，最小基可以通過有理運算導出，因此，特別是對於 $\Omega_\Gamma$，它只含有有理數係數。在兩個不定元的情形下，按照 Castelnuovo 的結果，仍可能涉及一個係數屬於代數數域 $\Omega$ 的基；實際研究（參見 F. Seidelmann 前引處）表明，對於每一個 $\Omega_\Gamma$，在這裡也能得到一個有理數係數基。幾乎無需計算便可從下述事實看出這一點：對於四元群，可以選取這樣的一個有理數係數基，即
\[
\varphi_1(x)=\sigma_1(x);\quad \varphi_2(x)=\frac{v w}{u};\quad \varphi_3(x)=\frac{w\cdot u}{v};\quad \varphi_4(x)=\frac{u\cdot v}{w},
\]
其中
\[
u=x_1+x_2-x_3-x_4;\quad v=x_1-x_2+x_3-x_4;\quad w=x_1-x_2-x_3+x_4,
\]
交錯群變為三個基函數 $\varphi_2,\varphi_3,\varphi_4$ 的循環群；每一個八階群則變為這幾個函數中任意兩個的對稱群，由此實現向三次和二次方程的群的約化。而對於循環群，一個有理數係數基早已由 Abel 給出，儘管當時並未這樣表述。\srcfn{*)}{參見 Weber Algebra（小版），§ 93，公式 (12)。} 因而證明了：\emph{具有給定群的三次和四次方程的全體 -- 其中至多須排除相對於參數表示為奇異的方程 -- 對於任意有理性域，都可以通過參數表示有理地構造；它們的流形與無附加約束方程的流形相同。}

實際研究表明（參見 F. Seidelmann 前引處），只有對於交錯群，並且僅當有理性域包含三次單位根域時，才需要一種補充表示；在所有其他情形下，參數表示則直接給出全部方程。

還應指出，所有\emph{阿貝爾群}的最小基也都是已知的。\srcfn{**)}{E. Fischer: Die Isomorphie der Invariantenkörper der endlichen Abelschen Gruppen linearer Transformationen. Gött. Nachr. 1915, 以及 Zur Theorie der endlichen Abelschen Gruppen. Math. Ann. 77 (1915).} 但這裡所涉及的是一個係數屬於由群特徵標導出的分圓域的基。因此，在這種情形下，對於構造具有給定群的方程，不能得到新的結果。

哥廷根，1916 年 7 月。

\clearpage
\setcounter{footnote}{0}
\providecommand{\bdelta}{\bar{\delta}}
\end{document}
